% =====================================================================
\chapter{Знакомство с пультами}
% =====================================================================
\chapterintro
  {из чего состоит пульт, как называются органы управления, чем отличаются TX12 и Boxer,
   какие бывают версии радиомодулей и как обращаться с аккумуляторами}
  {пульт TX12 или Boxer, заряженные аккумуляторы, кабель USB-C}

\section{Что такое аппаратура управления}

Пульт (передатчик, по-английски \emph{transmitter}, TX или \emph{radio}) измеряет положение
стиков, переключателей и крутилок, по заданным вами правилам превращает их в набор
\textbf{каналов} и по радио отправляет эти каналы \textbf{приёмнику} (\emph{receiver}, RX)
на модели. Приёмник выдаёт каналы на сервомашинки, регулятор оборотов или полётный контроллер.
Обратно от приёмника приходит \textbf{телеметрия}: качество связи, напряжение батареи модели,
данные GPS и т.\,д.

\begin{figure}[H]
\centering
\begin{tikzpicture}[node distance=6mm and 7mm]
  \node[blk,fill=blksw,minimum width=24mm] (ctl) {Стики,\\переключатели,\\крутилки};
  \node[blk,fill=blkmix,right=of ctl,minimum width=26mm] (etx) {EdgeTX:\\входы, миксеры,\\выходы};
  \node[blk,fill=blkrf,right=of etx,minimum width=22mm] (tx) {Радиомодуль\\(ELRS, Multi…)};
  \node[blk,fill=blkrf,right=13mm of tx,minimum width=16mm] (rx) {Приёмник};
  \node[blk,fill=blkout,right=of rx,minimum width=20mm] (srv) {Сервы, ESC,\\полётный\\контроллер};
  \draw[arr] (ctl) -- (etx);
  \draw[arr] (etx) -- node[above,lbl]{каналы} (tx);
  \draw[arr,decorate,decoration={snake,amplitude=1.2pt,segment length=5pt,pre length=2pt,post length=3pt}] ([yshift=2mm]tx.east) -- ([yshift=2mm]rx.west);
  \draw[arr,dashed] ([yshift=-2mm]rx.west) -- node[below,lbl]{телеметрия} ([yshift=-2mm]tx.east);
  \draw[arr] (rx) -- (srv);
  \begin{scope}[on background layer]
    \node[fit=(ctl)(etx)(tx),draw=etx,dashed,rounded corners,inner sep=6pt,label={[lbl,text=etx]above:Пульт}] {};
  \end{scope}
\end{tikzpicture}
\caption{Путь сигнала от руки пилота до модели}
\end{figure}

Современные пульты вроде TX12 и Boxer --- это небольшие компьютеры на микроконтроллере STM32
с монохромным экраном, SD-картой, динамиком и вибромотором. Всё «умное» в них делает
прошивка EdgeTX, а радиочасть вынесена в отдельный \textbf{радиомодуль}: внутренний
(встроенный) или внешний (вставляется в отсек на задней стороне).

\section{RadioMaster TX12}

TX12 --- компактный пульт «игрового» формата, один из самых популярных первых пультов для FPV.
Существует в двух поколениях:

\begin{description}[font=\sffamily\bfseries\color{etx}]
  \item[TX12 (MKI, 2021)] процессор STM32F2, стики на потенциометрах,
        внутренний модуль --- 4-в-1 мультипротокол или CC2500; исходно шёл с OpenTX, сейчас
        прошивается на EdgeTX.
  \item[TX12 Mark II (MKII, 2022)] процессор STM32F407, стики на датчиках Холла,
        заводская прошивка EdgeTX, внутренний модуль --- ExpressLRS 2,4\,ГГц до \SI{250}{мВт}
        \emph{или} CC2500; питание от двух элементов 18650 (или от 2S LiPo подходящего размера),
        зарядка через USB-C.
\end{description}

\begin{figure}[H]
\centering
\begin{tikzpicture}[scale=0.95]
  \fill[black!12,rounded corners=16pt] (-4.6,-2.3) rectangle (4.6,1.9);
  \fill[black!8,rounded corners=10pt] (-4.3,-3.7) rectangle (-2.3,-1.6);
  \fill[black!8,rounded corners=10pt] (2.3,-3.7) rectangle (4.3,-1.6);
  \fill[lcdbg,draw=black!60,rounded corners=1pt] (-1.2,-0.5) rectangle (1.2,0.85);
  \node[font=\sffamily\tiny] at (0,0.17) {ЖК 128×64};
  \foreach \x/\n in {-3.0/L,3.0/R}{
    \draw[fill=black!25] (\x,-0.2) circle (1.0);
    \fill[black!70] (\x,-0.2) circle (0.28);
  }
  % верхние органы
  \foreach \x/\l in {-4.0/SE,-3.3/SB,3.3/SC,4.0/SF}{
    \draw[fill=gray!60] (\x-0.07,1.85) rectangle (\x+0.07,2.3);
    \node[font=\sffamily\scriptsize\bfseries,above] at (\x,2.3) {\l};
  }
  \foreach \x/\l in {-2.55/SA,2.55/SD}{
    \draw[fill=gray!40] (\x-0.07,1.85) rectangle (\x+0.07,2.15);
    \node[font=\sffamily\scriptsize\bfseries,above] at (\x,2.15) {\l};
  }
  \foreach \x/\l in {-1.7/S1,1.7/S2}{
    \draw[fill=gray!50] (\x,1.95) ellipse (0.33 and 0.13);
    \node[font=\sffamily\scriptsize\bfseries,above] at (\x,2.05) {\l};
  }
  % кнопки
  \foreach \y/\l in {0.6/SYS,0.2/MDL,-0.2/TELE}{
    \draw[fill=gray!30,rounded corners=1pt] (-1.85,\y-0.12) rectangle (-1.45,\y+0.12);
  }
  \foreach \y/\l in {0.6/PAGE<,0.2/PAGE>,-0.2/RTN}{
    \draw[fill=gray!30,rounded corners=1pt] (1.45,\y-0.12) rectangle (1.85,\y+0.12);
  }
  \node[font=\sffamily\tiny,align=right,left] at (-1.9,0.6) {};
  \draw[fill=gray!40] (0,-1.1) circle (0.27);
  \draw[fill=gray!20] (0,-1.75) circle (0.18);
  % триммеры
  \foreach \x in {-3.0,3.0}{\draw[fill=gray!30,rounded corners=1pt] (\x-0.55,-1.5) rectangle (\x+0.55,-1.38);}
  \foreach \x in {-1.75,1.75}{\draw[fill=gray!30,rounded corners=1pt] (\x-0.06,-1.3) rectangle (\x+0.06,-0.55);}
  % подписи
  \draw[gray] (-1.85,0.6) -- (-5.4,0.9) node[left,lbl]{\btn{SYS}};
  \draw[gray] (-1.85,0.2) -- (-5.4,0.3) node[left,lbl]{\btn{MDL}};
  \draw[gray] (-1.85,-0.2) -- (-5.4,-0.3) node[left,lbl]{\btn{TELE}};
  \draw[gray] (1.85,0.6) -- (5.4,0.9) node[right,lbl]{\btn{PAGE<}};
  \draw[gray] (1.85,0.2) -- (5.4,0.3) node[right,lbl]{\btn{PAGE>}};
  \draw[gray] (1.85,-0.2) -- (5.4,-0.3) node[right,lbl]{\btn{RTN}};
  \draw[gray] (0.27,-1.1) -- (5.4,-1.0) node[right,lbl]{колесо + \btn{ENTER}};
  \draw[gray] (0.18,-1.75) -- (5.4,-1.8) node[right,lbl]{кнопка питания};
  \draw[gray] (-3.0,-1.44) -- (-5.4,-1.8) node[left,lbl]{триммеры T1…T4};
  \node[lbl] at (-3.0,-2.9) {левый стик};
  \node[lbl] at (3.0,-2.9) {правый стик};
\end{tikzpicture}
\caption{TX12: органы управления (схематично; расположение кнопок у ревизий может немного отличаться)}
\end{figure}

\begin{center}\small
\begin{tabular}{@{}lll@{}}
\toprule
Орган & Тип & Типичное назначение \\
\midrule
\sw{SA}, \sw{SD} & 2 положения (у части партий --- без фиксации) & арм, отсечка газа, запуск таймера \\
\sw{SB}, \sw{SC} & 3 положения & полётные режимы, расходы \\
\sw{SE}, \sw{SF} & 3 положения & режимы, закрылки, доп.\ каналы \\
\sw{S1}, \sw{S2} & потенциометры (крутилки) & настройка «на лету», камера, свет \\
Триммеры & 4 качельки (\sw{T1}--\sw{T4}) & подстройка нейтрали осей \\
\bottomrule
\end{tabular}
\end{center}

\begin{tip}
Точный тип каждого переключателя на \emph{вашем} экземпляре видно в
\mpath{SYS\sep Hardware}. Если вы поменяли тумблер (RadioMaster продаёт сменные двух- и
трёхпозиционные), там же нужно указать новый тип --- иначе EdgeTX неверно прочитает положения.
\end{tip}

\section{RadioMaster Boxer}

Boxer (2023) --- пульт покрупнее, с полноразмерными стиками на датчиках Холла (V4.0)
и большим аккумуляторным отсеком. Экран тот же, монохромный 128×64, поэтому интерфейс
EdgeTX на Boxer практически совпадает с TX12: всё, что сказано в книге про меню, относится
к обоим пультам.

Особенности Boxer:
\begin{itemize}
  \item внутренний модуль: \textbf{ExpressLRS до 1\,Вт} с активным охлаждением (вентилятором),
        либо 4-в-1 мультипротокол, либо CC2500 --- в зависимости от версии;
  \item питание: 2S LiPo (7,4\,В; в отсек помещается аккумулятор порядка \SI{6200}{мА\cdot ч})
        или два элемента 18650 в комплектном держателе;
  \item переключатели \sw{SA} (2 поз.), \sw{SB}, \sw{SC} (3 поз.), \sw{SD} (2 поз.);
        низкопрофильные кнопки \sw{SE} (с фиксацией) и \sw{SF} (без фиксации);
        два колеса-потенциометра \sw{S1}, \sw{S2};
  \item \textbf{шесть кнопок} под экраном, которые вместе работают как шестипозиционный
        переключатель (источник \sw{6POS}) --- удобно для выбора до шести полётных режимов;
  \item отсек для внешнего модуля формата JR.
\end{itemize}

\begin{figure}[H]
\centering
\begin{tikzpicture}[scale=0.95]
  \fill[black!12,rounded corners=20pt] (-5.0,-2.4) rectangle (5.0,2.0);
  \fill[black!8,rounded corners=12pt] (-4.8,-4.0) rectangle (-2.5,-1.6);
  \fill[black!8,rounded corners=12pt] (2.5,-4.0) rectangle (4.8,-1.6);
  \fill[lcdbg,draw=black!60,rounded corners=1pt] (-1.2,0.0) rectangle (1.2,1.4);
  \node[font=\sffamily\tiny] at (0,0.7) {ЖК 128×64};
  \foreach \i in {0,...,5}{
    \draw[fill=gray!30,rounded corners=1pt] (-1.15+\i*0.42,-0.45) rectangle ++(0.32,0.25);
    \node[font=\sffamily\tiny] at (-0.99+\i*0.42,-0.62) {\the\numexpr\i+1\relax};
  }
  \foreach \x in {-3.3,3.3}{
    \draw[fill=black!25] (\x,-0.3) circle (1.15);
    \fill[black!70] (\x,-0.3) circle (0.3);
  }
  \foreach \x/\l in {-4.3/SA,-3.6/SB,3.6/SC,4.3/SD}{
    \draw[fill=gray!60] (\x-0.07,1.95) rectangle (\x+0.07,2.4);
    \node[font=\sffamily\scriptsize\bfseries,above] at (\x,2.4) {\l};
  }
  \foreach \x/\l in {-2.6/SE,2.6/SF}{
    \draw[fill=gray!45,rounded corners=1pt] (\x-0.2,1.9) rectangle (\x+0.2,2.08);
    \node[font=\sffamily\scriptsize\bfseries,above] at (\x,2.08) {\l};
  }
  \foreach \x/\l in {-1.8/S1,1.8/S2}{
    \draw[fill=gray!50] (\x,2.0) ellipse (0.3 and 0.12);
    \node[font=\sffamily\scriptsize\bfseries,above] at (\x,2.08) {\l};
  }
  \draw[fill=gray!40] (0,-1.2) circle (0.27);
  \node[lbl] at (-3.3,-3.0) {левый стик};
  \node[lbl] at (3.3,-3.0) {правый стик};
  \draw[gray] (1.1,-0.35) -- (5.9,-0.1) node[right,lbl,align=left]{кнопки 1--6\\(\sw{6POS})};
  \draw[gray] (0.27,-1.2) -- (5.9,-1.3) node[right,lbl]{колесо + \btn{ENTER}};
  \draw[gray] (-2.6,2.0) -- (-5.9,1.3) node[left,lbl,align=right]{кнопка с фиксацией};
  \draw[gray] (2.6,2.0) -- (5.9,1.3) node[right,lbl,align=left]{кнопка без\\фиксации};
\end{tikzpicture}
\caption{Boxer: органы управления (схематично). Кнопки SYS, MDL, TELE, PAGE, RTN
расположены по бокам от экрана}
\end{figure}

\section{Сравнение TX12 MKII и Boxer}

\begin{center}\small
\renewcommand{\arraystretch}{1.2}
\begin{tabularx}{\linewidth}{@{}p{35mm}XX@{}}
\toprule
& \textbf{TX12 MKII} & \textbf{Boxer} \\
\midrule
Прошивка & EdgeTX & EdgeTX \\
Экран & монохромный 128×64 & монохромный 128×64 \\
Стики & Холл, компактные & Холл V4.0, полноразмерные \\
Переключатели & 6 (4 трёхпозиционных) & 4 тумблера + 2 кнопки \\
Крутилки & 2 & 2 колеса \\
Шестипозиционный & --- & да, 6 кнопок \\
Внутренний ELRS & до 250\,мВт & до 1\,Вт, вентилятор \\
Внешний модуль & отсек JR & отсек JR \\
Питание & 2×18650 & 2S LiPo или 2×18650 \\
Габариты & компактный, «геймпадный» хват & крупнее, удобен для «щипкового» хвата \\
\bottomrule
\end{tabularx}
\end{center}

\section{Аккумуляторы и зарядка}

\begin{itemize}
  \item \textbf{18650}: используйте качественные элементы \emph{одного} типа и одинаковой степени
        заряда. Соблюдайте полярность --- она выштампована в отсеке. Незащищённые элементы
        с плоским «плюсом» могут не доставать до контактов; с защитой --- могут не влезть по длине.
  \item \textbf{2S LiPo} (Boxer, частично TX12 MKII): проверьте разъём (часто XT30 или JST)
        и габариты. Раздутый аккумулятор в пульт ставить нельзя.
  \item Зарядка через USB-C встроенным зарядным устройством работает при \emph{выключенном} пульте.
        Индикатор заряда гаснет или меняет цвет, когда заряд окончен (зависит от ревизии).
\end{itemize}

\begin{caution}
Не оставляйте заряжающийся пульт без присмотра и не заряжайте повреждённые элементы.
Если пульт долго не используется, храните литиевые аккумуляторы заряженными примерно
до \SI{3.8}{В} на элемент (7,6\,В на 2S).
\end{caution}

\section{Режимы стиков (Mode 1/2)}

Какой стик управляет каким каналом, определяет \textbf{режим} (Mode):

\begin{center}\small
\begin{tabular}{@{}lcccc@{}}
\toprule
& Левый ↕ & Левый ↔ & Правый ↕ & Правый ↔\\
\midrule
Mode 1 & тангаж (Ele) & рыскание (Rud) & газ (Thr) & крен (Ail)\\
\textbf{Mode 2} & \textbf{газ (Thr)} & \textbf{рыскание (Rud)} & \textbf{тангаж (Ele)} & \textbf{крен (Ail)}\\
Mode 3 & тангаж & крен & газ & рыскание\\
Mode 4 & газ & крен & тангаж & рыскание\\
\bottomrule
\end{tabular}
\end{center}

Mode 2 --- самый распространённый в мире и в FPV. Механически стик газа у Mode 2 находится
слева и не возвращается в центр (есть трещотка или фрикцион). Если вы купили пульт в Mode 2,
а хотите Mode 1, придётся перенести пружину центрирования и трещотку на другой стик
(инструкции есть у производителя) \emph{и} поменять режим в EdgeTX
(\mpath{SYS\sep Radio Setup\sep Mode}, глава~\ref{ch:radio}).

\begin{summary}
\begin{itemize}
  \item Пульт = органы управления + EdgeTX + радиомодуль; приёмник на модели получает каналы
        и возвращает телеметрию.
  \item TX12 и Boxer имеют одинаковый экран, поэтому меню EdgeTX у них одинаковые.
  \item Главные отличия --- набор переключателей, 6POS-кнопки и мощность ELRS у Boxer.
  \item Режим стиков (обычно Mode 2) задаётся и механически, и в прошивке.
\end{itemize}
\end{summary}
